% !TEX program = xelatex
\documentclass[11pt,a4paper]{article}

\usepackage[UTF8]{ctex}
\usepackage[margin=2.2cm]{geometry}
\usepackage{amsmath, amssymb, amsthm}
\usepackage{listings}
\usepackage{xcolor}
\usepackage{tikz}
\usepackage{booktabs}
\usepackage{multirow}
\usepackage{hyperref}
\usetikzlibrary{positioning, arrows.meta, shapes.geometric, fit, calc, decorations.pathreplacing}

\hypersetup{
  colorlinks=true,
  linkcolor=blue!60!black,
  urlcolor=teal!70!black,
}

\lstdefinelanguage{Rust}{
  morekeywords={fn,let,mut,pub,struct,impl,use,mod,self,Self,Option,Some,None,return,if,else,for,while,loop,match,break,continue,as,where,trait,type,enum,move,ref,const,static,unsafe,extern,crate,super,dyn,async,await,box,in},
  sensitive=true,
  morecomment=[l]{//},
  morecomment=[s]{/*}{*/},
  morestring=[b]",
}

\lstset{
  language=Rust,
  basicstyle=\ttfamily\footnotesize,
  keywordstyle=\color{blue!60!black}\bfseries,
  commentstyle=\color{green!40!black}\itshape,
  stringstyle=\color{red!60!black},
  numbers=left,
  numberstyle=\tiny\color{gray},
  numbersep=8pt,
  frame=single,
  framerule=0.4pt,
  rulecolor=\color{gray!40},
  breaklines=true,
  breakatwhitespace=true,
  showstringspaces=false,
  tabsize=4,
  columns=fullflexible,
  keepspaces=true,
}

\title{\textbf{halfedge\,: 测试夹具模块设计文档} \\ \large \texttt{src/test\_util.rs}}
\author{halfedge 库}
\date{2026-07-04}

\begin{document}
\maketitle
\tableofcontents

\section{模块定位}

\texttt{test\_util.rs} 提供 \texttt{build\_icosphere}：生成单位 icosphere
测试网格，用于快速调试几何算法、IO、校验、导出等功能。

\begin{center}
\renewcommand{\arraystretch}{1.18}
\begin{tabular}{@{}ll@{}}
  \toprule
  \textbf{API} & \texttt{build\_icosphere(subdivisions: usize) -> MeshStorage} \\
  \midrule
  球心 / 半径 & 原点 / $1$ \\
  基础形状     & icosahedron（12 顶点 20 面） \\
  细分策略     & 每面分裂为 4，中点归一化到单位球面 \\
  用途         & 测试与调试（非生产场景） \\
  \bottomrule
\end{tabular}
\end{center}

\section{icosahedron 构造}

\subsection{黄金比顶点}

icosahedron 的 12 个顶点由黄金比 $\varphi = \frac{1 + \sqrt{5}}{2}$ 构造，
分布在三个相互正交的矩形上：

\begin{center}
\begin{tikzpicture}[
  vert/.style={circle, draw=blue!70!black, fill=blue!10, minimum size=5mm, inner sep=0pt, font=\tiny},
  axis/.style={->, thick, draw=gray!60!black}
]
  % 三个正交矩形
  \draw[gray!50] (0,0) ellipse (2 and 1.2);
  \draw[gray!50, dashed] (0,0) ellipse (1.2 and 2);
  \draw[gray!50, dotted] (-2,0) -- (2,0);
  \draw[gray!50, dotted] (0,-2) -- (0,2);

  % 12 顶点（投影示意）
  \node[vert] at (2,0)    {$v_0$};
  \node[vert] at (-2,0)   {$v_1$};
  \node[vert] at (1.2,1.2){$v_2$};
  \node[vert] at (-1.2,1.2){$v_3$};
  \node[vert] at (1.2,-1.2){$v_4$};
  \node[vert] at (-1.2,-1.2){$v_5$};
  \node[vert] at (0,2)    {$v_6$};
  \node[vert] at (0,-2)   {$v_7$};
  \node[vert] at (0.7,0.7){$v_8$};
  \node[vert] at (-0.7,0.7){$v_9$};
  \node[vert] at (0.7,-0.7){$v_{10}$};
  \node[vert] at (-0.7,-0.7){$v_{11}$};

  \node[below=2.5cm of current bounding box.south, font=\small, align=center]
    {黄金比 $\varphi \approx 1.618$，顶点形如 $(\pm 1, \pm \varphi, 0)$ 的循环置换};
\end{tikzpicture}
\end{center}

代码中 \texttt{ICOSA\_VERTICES} 数组（未归一化）：
\begin{lstlisting}
const ICOSA_VERTICES: [[f64; 3]; 12] = [
    [-1.0,  PHI, 0.0],  [ 1.0,  PHI, 0.0],
    [-1.0, -PHI, 0.0],  [ 1.0, -PHI, 0.0],
    [ 0.0, -1.0,  PHI],  [ 0.0,  1.0,  PHI],
    [ 0.0, -1.0, -PHI],  [ 0.0,  1.0, -PHI],
    [ PHI, 0.0, -1.0],  [ PHI, 0.0,  1.0],
    [-PHI, 0.0, -1.0],  [-PHI, 0.0,  1.0],
];
\end{lstlisting}

构造时所有顶点被归一化到单位球面：
\[
  \vec{p}' = \frac{\vec{p}}{\|\vec{p}\|}, \qquad
  \|\vec{p}\| = \sqrt{1 + \varphi^2} = \sqrt{\frac{5 + \sqrt{5}}{2}}.
\]

\subsection{20 个三角面}

\texttt{ICOSA\_FACES} 列出 20 个 CCW 朝外的三角面。CCW 朝向保证：
\begin{itemize}
  \item 面法向 $\hat{n} = (\vec{B}-\vec{A}) \times (\vec{C}-\vec{A})$ 指向外侧；
  \item 每条无向边在两个相邻面中方向相反（绕向一致）。
\end{itemize}

\section{细分算法}

\subsection{单次细分：1 面 $\to$ 4 面}

对每个三角面 $[A, B, C]$：
\begin{enumerate}
  \item 计算三条边中点 $M_{AB}, M_{BC}, M_{CA}$；
  \item 每个中点归一化到单位球面：$\vec{M}' = \vec{M} / \|\vec{M}\|$；
  \item 用 4 个小三角形替换原三角形：
    \begin{itemize}
      \item $[A, M_{AB}, M_{CA}]$
      \item $[B, M_{BC}, M_{AB}]$
      \item $[C, M_{CA}, M_{BC}]$
      \item $[M_{AB}, M_{BC}, M_{CA}]$（中心）
    \end{itemize}
\end{enumerate}

\begin{center}
\begin{tikzpicture}[
  vert/.style={circle, draw=blue!70!black, fill=blue!10, minimum size=5mm, inner sep=0pt, font=\footnotesize},
  newvert/.style={circle, draw=red!70!black, fill=red!10, minimum size=5mm, inner sep=0pt, font=\footnotesize},
  he/.style={thick, draw=gray!60!black}
]
  % 左：细分前
  \node[vert] (a1) at (0,0) {$A$};
  \node[vert] (b1) at (3,0) {$B$};
  \node[vert] (c1) at (1.5,2.5) {$C$};
  \draw[he] (a1) -- (b1) -- (c1) -- (a1);
  \node[below=0.3cm of a1, xshift=1.5cm, font=\small] {细分前：1 面};

  % 右：细分后
  \node[vert] (a2) at (5,0) {$A$};
  \node[vert] (b2) at (8,0) {$B$};
  \node[vert] (c2) at (6.5,2.5) {$C$};
  \node[newvert] (mab) at (6.5,0) {$M_{AB}$};
  \node[newvert] (mbc) at (7.25,1.25) {$M_{BC}$};
  \node[newvert] (mca) at (5.75,1.25) {$M_{CA}$};
  \draw[he] (a2) -- (mab) -- (b2);
  \draw[he] (b2) -- (mbc) -- (c2);
  \draw[he] (c2) -- (mca) -- (a2);
  \draw[he, red!70!black] (mab) -- (mbc) -- (mca) -- (mab);
  \node[below=0.3cm of a2, xshift=1.5cm, font=\small] {细分后：4 面};
\end{tikzpicture}
\end{center}

\subsection{中点缓存}

同一条无向边被两个相邻面共享，若不缓存会产生两个不同的中点顶点，
破坏拓扑一致性。\texttt{midpoint} 函数用 \texttt{HashMap<(u32, u32), u32>}
缓存「端点对 $\to$ 中点索引」，键按端点 ID 大小排序（无向边）。

\subsection{细分次数与网格规模}

设细分次数为 $n$，则
\begin{align}
  V_n &= 10 \cdot 4^n + 2, \\
  F_n &= 20 \cdot 4^n, \\
  E_n &= 30 \cdot 4^n.
\end{align}

\begin{center}
\renewcommand{\arraystretch}{1.18}
\begin{tabular}{@{}rrrrr@{}}
  \toprule
  $n$ & $V_n$ & $F_n$ & $E_n$ & 半边数 \\
  \midrule
  0 & 12   & 20   & 30   & 60    \\
  1 & 42   & 80   & 120  & 240   \\
  2 & 162  & 320  & 480  & 960   \\
  3 & 642  & 1280 & 1920 & 3840  \\
  4 & 2562 & 5120 & 7680 & 15360 \\
  \bottomrule
\end{tabular}
\end{center}

\textbf{公式推导}：每次细分，面数 $F \to 4F$；边数 $E \to 2E$（每条原边分裂为 2）；
顶点数 $V \to V + E$（每条原边产生 1 个新中点）。递推解得上述闭式。

\section{整体流程}

\begin{center}
\resizebox{\textwidth}{!}{%
\begin{tikzpicture}[
  node distance=0.5cm,
  proc/.style={draw, rounded corners=2pt, minimum width=10cm, minimum height=0.7cm, align=center, font=\small},
  dec/.style={diamond, draw, aspect=2.4, fill=orange!12, inner sep=1pt, font=\small, align=center, minimum width=3.5cm},
  op/.style={-Stealth, thick},
  startend/.style={draw, rounded corners=10pt, minimum width=2.6cm, minimum height=0.7cm, font=\small, fill=gray!12}
]
  \node[startend] (s) {开始：$n$};
  \node[proc, below=of s, fill=blue!8] (p1) {初始化 \texttt{vertices = ICOSA\_VERTICES}（归一化），\texttt{faces = ICOSA\_FACES}};
  \node[dec, below=of p1] (d1) {剩余细分 $> 0$？};
  \node[proc, below=of d1, fill=yellow!15] (p2) {\texttt{subdivide\_once}：每面 $\to$ 4 面，中点归一化};
  \node[proc, below=of p2, fill=yellow!15] (p3) {细分计数 $-1$；回到 $d1$};
  \node[proc, right=1.6cm of d1, fill=green!12] (done) {调用 \texttt{build\_mesh\_from\_vertices\_and\_faces}};
  \node[startend, below=of done] (e) {返回 \texttt{MeshStorage}};
  \draw[op] (s) -- (p1);
  \draw[op] (p1) -- (d1);
  \draw[op] (d1) -- node[right, font=\scriptsize] {是} (p2);
  \draw[op] (p2) -- (p3);
  \draw[op, dashed] (p3.east) -- ++(0.6,0) |- (d1.east);
  \draw[op] (d1) -- node[above, font=\scriptsize] {否} (done);
  \draw[op] (done) -- (e);
\end{tikzpicture}
}
\end{center}

\section{复杂度}

顶点数 $O(4^n)$，构建时间 $O(4^n)$。

\section{退化守卫}

\begin{itemize}
  \item \texttt{normalize} 函数：若输入长度 $< 10^{-12}$，返回原向量（避免除零）；
  \item \texttt{midpoint} 缓存：同一边的两次访问返回同一中点索引，避免重复创建；
  \item \texttt{ICOSA\_FACES} 绕向经过测试验证一致（见
        \texttt{icosphere\_face\_normals\_outward}）。
\end{itemize}

\section{测试覆盖}

\begin{center}
\renewcommand{\arraystretch}{1.18}
\begin{tabular}{@{}ll@{}}
  \toprule
  \textbf{测试名} & \textbf{验证点} \\
  \midrule
  \texttt{icosphere\_subdiv0\_basic}            & 12 顶点 / 20 面 / 60 半边 \\
  \texttt{icosphere\_subdiv0\_passes\_validation} & 0 细分通过完整校验 \\
  \texttt{icosphere\_subdiv1\_basic}            & 42 顶点 / 80 面 / 240 半边 \\
  \texttt{icosphere\_subdiv1\_passes\_validation} & 1 细分通过完整校验 \\
  \texttt{icosphere\_subdiv2\_passes\_validation} & 2 细分通过完整校验 \\
  \texttt{icosphere\_vertices\_on\_unit\_sphere} & 所有顶点 $|\vec{p}|^2 = 1$ \\
  \texttt{icosphere\_face\_normals\_outward}    & 所有面法向与重心方向同向 \\
  \texttt{icosphere\_export\_buffers\_consistent} & 导出缓冲规模正确（42 顶点 / 240 索引） \\
  \bottomrule
\end{tabular}
\end{center}

\end{document}
